\textbf{COLUNI 2019 -- Questão 26}\par\medskip O magnésio é o oitavo elemento químico mais abundante na crosta terrestre e é o terceiro elemento em maior quantidade na água do mar. A obtenção de íons magnésio, $\mathrm{Mg^{2+}}$, a partir da água do mar se baseia no fato de que o hidróxido de magnésio, $\mathrm{Mg(OH)_2}$, é menos solúvel em água do que o hidróxido de cálcio, $\mathrm{Ca(OH)_2}$.

\par

A adição de cal, $\mathrm{CaO}$, a uma amostra de água do mar provoca a precipitação de hidróxido de magnésio:

\begin{align*}
\text{Equação 1:}\quad & \mathrm{CaO(s) + H_2O(l) \longrightarrow Ca^{2+}(aq) + 2OH^{-}(aq)}\\
\text{Equação 2:}\quad & \mathrm{Mg^{2+}(aq) + 2OH^{-}(aq) \longrightarrow Mg(OH)_2(s)}
\end{align*}

\par

Posteriormente, o hidróxido de magnésio é convertido em cloreto de magnésio, $\mathrm{MgCl_2}$, pela adição de ácido clorídrico, HCl:

\par

\[\mathrm{Mg(OH)_2(s) + 2HCl(aq) \longrightarrow MgCl_2(aq) + 2H_2O(l)}\]

\par

Para obter o magnésio metálico, $\mathrm{Mg}$, o cloreto de magnésio é fundido e submetido à passagem de corrente elétrica, processo denominado eletrólise, no qual são obtidos simultaneamente o magnésio metálico e o cloro gasoso, $\mathrm{Cl_2}$:

\par

\[\mathrm{MgCl_2(l) \longrightarrow Mg(s) + Cl_2(g)}\]

\par

Com base nas informações fornecidas e nos conhecimentos sobre ciências, assinale a afirmativa CORRETA:\begin{enumerate}
\item O processo mais indicado para separar o hidróxido de magnésio precipitado da fase líquida é a destilação
      fracionada.
\item A obtenção de cloreto de magnésio a partir do hidróxido de magnésio e ácido clorídrico é uma reação de
      combustão.
\item O cloro sofre uma mudança de estado físico chamada de sublimação, durante o processo de obtenção do
      magnésio metálico.
\item A obtenção de magnésio metálico e do cloro gasoso a partir do cloreto de magnésio fundido é um
      fenômeno químico.\end{enumerate}
